%! Function Fundamentals: Notation, Domain, and Range — Unit 2, Lesson 2.1 — Warm-Up (Answer Key)
\documentclass[10pt]{article}
\usepackage{atda-article}
\usepackage{atda-key}   % pulls in -boxes; do NOT also load -boxes
\newcommand{\TallMath}[1]{$\displaystyle #1\rule[-1.4em]{0pt}{3.2em}$}
\setlength{\workrowsep}{6pt}   % handwriting room; moves blank and key together

\begin{document}

\pageheader{Unit 2, Lesson 2.1}{Warm-Up --- Answer Key}
% NAMESTRIP: no \namedateperiod here — mirrors the blank. Teacher-only prose goes
% in the lesson plan as \begin{teachernote}[Warm-Up], never in a key.

\vspace{0.15in}

\begin{notesbox}{Which Slot Does the Number Go In?}
\small
Five quick items. Items 1 and 4 are the same question asked two ways --- watch for that.

\begin{enumerate}[leftmargin=1.6em, itemsep=4pt, topsep=4pt]

\item \textbf{Both directions.} Let $f(x)=4x-9$. Find $f(5)$, then find the value of $x$ for which
$f(x)=7$.
\begin{work}
  f(5) &= 4(5)-9 \\
       &= 11 \\
  4x-9 &= 7 \\
    4x &= 16 \\
     x &= 4
\end{work}

\item \textbf{A value the algebra will not allow.} For which value of $x$ is
$\dfrac{2x}{x+6}$ undefined?
\begin{work}
  x+6 &= 0 \\
    x &= -6
\end{work}

\item \textbf{Writing a set of numbers.} Write an inequality for ``every number from $-2$ up to
$5$, with $-2$ included and $5$ left out.''

\ansline{$-2 \le x < 5$}\\

\item \textbf{Reading a table both ways.} Here is a function $g$ given by a table.
\smallskip

\renewcommand{\arraystretch}{1.25}
\begin{tabular}{@{}|c|c|c|c|c|c|@{}}
\hline
$x$ & $-2$ & $-1$ & $0$ & $1$ & $2$ \\ \hline
$g(x)$ & $5$ & $1$ & $-1$ & $1$ & $5$ \\ \hline
\end{tabular}
\smallskip

Find $g(-1)$. Then find \emph{every} value of $x$ for which $g(x)=5$.

\ansline{$g(-1)=1$. \quad $g(x)=5$ at both $x=-2$ and $x=2$ --- two inputs, one output.}\\

\item \textbf{In context.} A gym charges \$25 to join plus \$15 per month, so the total paid after
$m$ months is $T(m)=15m+25$. Find $T(6)$, then write one sentence saying what that number means
for a member.
\begin{work}
  T(6) &= 15(6)+25 \\
       &= 115
\end{work}
\ansline{After six months a member has paid \$115 in all --- the \$25 to join plus six \$15 months.}

\end{enumerate}
\end{notesbox}

\end{document}
